% Module 5b: Computing with Qubits -- NISQ Variational Algorithms
% Estimated slides: 30 | Duration: ~2 hours
% Key topics: NISQ overview, parameterized circuits, VQE, QAOA,
%             hybrid quantum-classical optimization, ansatz design
% Labs: QAOA for Max-Cut optimization; VQE for H2 molecule

%%%%%%%%%%%%%%%%%%%%%%%%%%%%%%%%%%%%%%%%%%%%%%%%%%%%%%%%%%
\begin{frame}[fragile]\frametitle{}
\begin{center}
{\Large Module 5b: NISQ Variational Algorithms}\\[0.5em]
{\large VQE, QAOA, and the Hybrid Quantum-Classical Approach}
\end{center}
\end{frame}

%%%%%%%%%%%%%%%%%%%%%%%%%%%%%%%%%%%%%%%%%%%%%%%%%%%%%%%%%%%
\begin{frame}[fragile]\frametitle{The NISQ Era: What We Have Today}
\begin{itemize}
\item NISQ = Noisy Intermediate-Scale Quantum (Preskill, 2018)
\item Hundreds to thousands of physical qubits, but error rates of 0.1-1\%
\item Circuit depth limited: errors accumulate beyond ~100 gate layers
\item No error correction yet (that requires millions of physical qubits)
\item Question: what useful computation can we do on noisy, shallow circuits?
\item Answer: variational (hybrid quantum-classical) algorithms
\end{itemize}
% Suggested diagram: NISQ limitations table (qubits, depth, fidelity, connectivity)
\end{frame}

%%%%%%%%%%%%%%%%%%%%%%%%%%%%%%%%%%%%%%%%%%%%%%%%%%%%%%%%%%%
\begin{frame}[fragile]\frametitle{Hybrid Quantum-Classical Computing}
\begin{itemize}
\item Idea: use quantum for the hard part, classical optimizer for the rest
\item Quantum circuit has tunable parameters $\vec{\theta}$ (rotation angles)
\item Classical optimizer adjusts $\vec{\theta}$ to minimize a cost function
\item Quantum computer evaluates the cost function efficiently; classical minimizes it
\item Similar to training a neural network: quantum circuit = forward pass, optimizer = backprop
\item Examples: VQE, QAOA, Quantum Neural Networks
\end{itemize}
% Suggested diagram: Hybrid loop diagram: parametric circuit -> measure -> cost -> optimizer -> update theta
\end{frame}

%%%%%%%%%%%%%%%%%%%%%%%%%%%%%%%%%%%%%%%%%%%%%%%%%%%%%%%%%%%
\begin{frame}[fragile]\frametitle{Parameterized Quantum Circuits (PQC)}
\begin{itemize}
\item A quantum circuit with rotation angles as free parameters: $R_y(\theta_1)$, $R_z(\theta_2)$, etc.
\item The circuit defines a family of quantum states parameterized by $\vec{\theta}$
\item Changing $\vec{\theta}$ changes the output state (and thus measurement probabilities)
\item PQCs are the core building block of all variational algorithms
\item Qiskit: use \texttt{ParameterVector} or \texttt{Parameter} objects
\item Hardware efficient: layers of rotation gates + CNOT gates repeat $L$ times
\end{itemize}
\begin{verbatim}
from qiskit.circuit import ParameterVector
theta = ParameterVector('theta', 4)
qc.ry(theta[0], 0); qc.ry(theta[1], 1)
qc.cx(0, 1)
\end{verbatim}
% Suggested diagram: Parameterized circuit with theta parameters labeled
\end{frame}

%%%%%%%%%%%%%%%%%%%%%%%%%%%%%%%%%%%%%%%%%%%%%%%%%%%%%%%%%%%
\begin{frame}[fragile]\frametitle{Variational Quantum Eigensolver (VQE)}
\begin{itemize}
\item Goal: find the ground state energy of a molecule or material
\item Why quantum? Simulating quantum systems on classical computers scales exponentially
\item VQE: prepare a parameterized trial state (ansatz), measure its energy, minimize
\item Energy = expectation value of the Hamiltonian (observable describing energy)
\item Variational principle: any trial state has energy $\geq$ true ground state energy
\item Minimize energy over $\vec{\theta}$ to approach the true ground state
\end{itemize}
% Suggested diagram: Variational principle energy landscape; optimizer converging to minimum
\end{frame}

%%%%%%%%%%%%%%%%%%%%%%%%%%%%%%%%%%%%%%%%%%%%%%%%%%%%%%%%%%%
\begin{frame}[fragile]\frametitle{VQE: The Algorithm Step by Step}
\begin{itemize}
\item Step 1: Express molecule as a Hamiltonian (sum of Pauli operators)
\item Step 2: Choose an ansatz circuit (e.g., UCCSD or hardware-efficient)
\item Step 3: Run quantum circuit with current $\vec{\theta}$, measure Pauli terms
\item Step 4: Compute expectation value $\langle H \rangle = \langle \psi(\theta) | H | \psi(\theta) \rangle$
\item Step 5: Classical optimizer updates $\vec{\theta}$ to reduce $\langle H \rangle$
\item Repeat until convergence; final $\langle H \rangle$ approximates ground state energy
\end{itemize}
% Suggested diagram: VQE flowchart with quantum and classical parts highlighted
\end{frame}

%%%%%%%%%%%%%%%%%%%%%%%%%%%%%%%%%%%%%%%%%%%%%%%%%%%%%%%%%%%
\begin{frame}[fragile]\frametitle{VQE: Applications}
\begin{itemize}
\item \textbf{Chemistry}: ground state energy of H$_2$, LiH, BeH$_2$ molecules
\item \textbf{Material science}: electronic structure of materials, superconductors
\item \textbf{Drug discovery}: binding energies of drug-target interactions
\item \textbf{Catalyst design}: reaction energies for industrial catalysts
\item Today: small molecules (up to ~30 qubits) demonstrated on real hardware
\item Future: larger molecules as qubit counts and fidelities improve
\end{itemize}
% Suggested diagram: H2 molecule diagram; energy vs bond length VQE vs exact comparison
\end{frame}

%%%%%%%%%%%%%%%%%%%%%%%%%%%%%%%%%%%%%%%%%%%%%%%%%%%%%%%%%%%
\begin{frame}[fragile]\frametitle{Quantum Approximate Optimization Algorithm (QAOA)}
\begin{itemize}
\item Goal: solve combinatorial optimization problems (e.g., find maximum cut in a graph)
\item QAOA: a specific parameterized circuit structure designed for optimization
\item Alternates between cost unitary $U_C$ (encodes problem) and mixer unitary $U_B$ (explores)
\item Circuit depth $p$ (number of alternating layers); higher $p$ = better solution
\item For $p \to \infty$: provably converges to optimal solution
\item In practice: even $p=1,2$ gives good approximate solutions on NISQ
\end{itemize}
% Suggested diagram: QAOA circuit showing alternating UC and UB layers
\end{frame}

%%%%%%%%%%%%%%%%%%%%%%%%%%%%%%%%%%%%%%%%%%%%%%%%%%%%%%%%%%%
\begin{frame}[fragile]\frametitle{QAOA: Max-Cut Example}
\begin{itemize}
\item Max-Cut: partition graph nodes into two groups to maximize the number of edges between groups
\item NP-hard classically; QAOA finds approximate solutions
\item Encode: each qubit = one node; $|0\rangle$ = group A, $|1\rangle$ = group B
\item Cost: maximize number of edges between qubits in different states
\item QAOA with $p=1$: guaranteed to find a cut of at least 0.6924 $\times$ optimal (proven)
\item Industry use: vehicle routing, portfolio optimization, job scheduling
\end{itemize}
% Suggested diagram: Graph with highlighted max-cut; qubit assignment visualization
\end{frame}

%%%%%%%%%%%%%%%%%%%%%%%%%%%%%%%%%%%%%%%%%%%%%%%%%%%%%%%%%%%
\begin{frame}[fragile]\frametitle{Ansatz Design}
\begin{itemize}
\item Ansatz = the parameterized circuit structure used in variational algorithms
\item Needs to balance expressibility (can reach target state) vs trainability (gradients exist)
\item Problem-inspired ansatz (UCCSD for chemistry): uses domain knowledge, deeper
\item Hardware-efficient ansatz: uses native gates, shallower, but less expressive
\item Barren plateaus: gradients vanish exponentially in $n$ qubits for random ansatze
\item Active research: designing ansatze that avoid barren plateaus
\end{itemize}
% Suggested diagram: Expressibility vs depth trade-off; hardware-efficient vs UCCSD comparison
\end{frame}

%%%%%%%%%%%%%%%%%%%%%%%%%%%%%%%%%%%%%%%%%%%%%%%%%%%%%%%%%%%
\begin{frame}[fragile]\frametitle{Classical Optimizers for Variational Algorithms}
\begin{itemize}
\item Gradient-based: BFGS, L-BFGS-B; requires gradient estimation (parameter shift rule)
\item Gradient-free: COBYLA, Nelder-Mead, SPSA; more robust to noise
\item Parameter shift rule: estimate gradient by shifting parameter by $\pm\pi/2$
\item SPSA (Simultaneous Perturbation Stochastic Approximation): efficient for noisy settings
\item Convergence: variational algorithms can get stuck in local minima
\item Initialization: random initialization can land in barren plateaus; warm-starting helps
\end{itemize}
% Suggested diagram: Loss landscape with local minima and optimization trajectories
\end{frame}

%%%%%%%%%%%%%%%%%%%%%%%%%%%%%%%%%%%%%%%%%%%%%%%%%%%%%%%%%%%
\begin{frame}[fragile]\frametitle{Lab: QAOA for Max-Cut in Qiskit}
\begin{itemize}
\item Use Qiskit's QAOA implementation on a 4-node graph
\item Encode the Max-Cut cost Hamiltonian
\item Optimize with COBYLA; compare result to brute-force classical
\item Visualize the final cut assignment
\end{itemize}
\begin{verbatim}
from qiskit_algorithms import QAOA
from qiskit_optimization.applications import Maxcut
import networkx as nx
G = nx.Graph(); G.add_edges_from([(0,1),(1,2),(2,3),(3,0),(0,2)])
\end{verbatim}
\end{frame}

%%%%%%%%%%%%%%%%%%%%%%%%%%%%%%%%%%%%%%%%%%%%%%%%%%%%%%%%%%%
\begin{frame}[fragile]\frametitle{Limitations of NISQ Variational Algorithms}
\begin{itemize}
\item Barren plateaus: gradient vanishes exponentially with qubit count
\item Classical simulation can sometimes beat QAOA at low depth $p$
\item Quantum advantage for variational algorithms not yet proven rigorously
\item Measurement overhead: many Pauli term measurements needed (for chemistry)
\item Noise increases with depth; more layers $p$ = more errors
\item These are open research problems --- active area of innovation
\end{itemize}
\end{frame}

%%%%%%%%%%%%%%%%%%%%%%%%%%%%%%%%%%%%%%%%%%%%%%%%%%%%%%%%%%%
\begin{frame}[fragile]\frametitle{Module 5b Summary}
\begin{itemize}
\item NISQ era: noisy qubits, shallow circuits, no error correction yet
\item Variational algorithms: hybrid quantum-classical loop with parameterized circuits
\item VQE: find molecular ground state energies; chemical simulation
\item QAOA: combinatorial optimization (Max-Cut, routing, scheduling)
\item Parameterized circuits + classical optimizer = quantum machine learning paradigm
\item Ansatz design and optimizer choice are critical for performance
\item Limitations: barren plateaus, measurement overhead, noise
\item Next: Quantum Fourier Transform and Shor's Algorithm
\end{itemize}
\end{frame}
